% ============================================================
% Project Status Summary -- Text Signals, Outcome Fixes, Prediction Models
% Prepared 2026-09-17
% ============================================================
\documentclass[aspectratio=169,10pt]{beamer}

% -- Theme -------------------------------------------------------------------
\usetheme{Madrid}
\usecolortheme{default}
\setbeamertemplate{navigation symbols}{}
\setbeamertemplate{footline}[frame number]
\setbeamersize{text margin left=6mm, text margin right=6mm}

% -- Packages ------------------------------------------------------------------
\usepackage{amsmath, amssymb}
\usepackage{graphicx}
\usepackage{booktabs}
\usepackage{array}
\usepackage{tabularx}
\usepackage{xcolor}
\usepackage{colortbl}
\usepackage{hyperref}
\pdfstringdefDisableCommands{\let\\\space\let$\relax}

% -- Colors -- matches Paper/Text Analysis/20260714_informs_slides.tex --------
\definecolor{ncsuRed}{RGB}{204,0,0}
\definecolor{ncsuGray}{RGB}{100,100,100}
\definecolor{techBlue}{RGB}{37,99,235}
\definecolor{govGreen}{RGB}{5,150,105}
\definecolor{warnOrange}{RGB}{217,119,6}
\definecolor{lightGray}{RGB}{243,244,246}

\setbeamercolor{structure}{fg=ncsuRed}
\setbeamercolor{block title}{bg=ncsuRed!90, fg=white}
\setbeamercolor{block body}{bg=ncsuRed!8}
\setbeamercolor{alerted text}{fg=ncsuRed}

\newcommand{\highlight}[1]{{\color{ncsuRed}\textbf{#1}}}
\newcommand{\tech}[1]{{\color{techBlue}\textbf{#1}}}
\newcommand{\gov}[1]{{\color{govGreen}\textbf{#1}}}

% -- Title ---------------------------------------------------------------------
\title[Project Status Summary]{Project Status Summary}
\subtitle{Text extraction, the outcome-variable problem, and prediction models}

\author[Sahebi-Fakhrabad]{Amirreza Sahebi-Fakhrabad}
\date{September 17, 2026 -- internal, not for external distribution}

\begin{document}

% ─────────────────────────────────────────────────────────────────────────────
\begin{frame}
\titlepage
\end{frame}

% ─────────────────────────────────────────────────────────────────────────────
\begin{frame}{Roadmap}
\begin{enumerate}
  \setlength{\itemsep}{12pt}
  \item \textbf{The extraction pipeline} -- validated, model chosen, \highlight{closed}
  \item \textbf{Is the text useful?} -- and a bigger problem we found in the outcome variable
  \item \textbf{Prediction models} -- recall, adverse events, shortage
  \item \textbf{Where this points next}
\end{enumerate}
\end{frame}

% ═══════════════════════════════════════════════════════════════════════════
\section{Part 1: Extraction Pipeline -- Closed}

\begin{frame}{Part 1: The Extraction Pipeline is Closed}
\begin{itemize}
  \setlength{\itemsep}{10pt}
  \item \textbf{Scale:} 1,067 observations, 98 facilities, 14 APIs, two LLM providers
  \item \textbf{Validation:} independent human reviewer blind-labeled 50 observations
  \item \textbf{Found and fixed:} 1 real bug, 6 real prompt gaps (contamination rules,
        patient-risk scoping, scope calibration, root-cause bleed, data-integrity gap)
  \item \textbf{Found and standardized:} severity\_tier's real confusion is
        Major/Moderate, not Critical/Major -- collapsing the right pair takes accuracy
        from 66--68\% to \highlight{90--94\%}
  \item \textbf{Model decision:} \tech{Claude Sonnet 5} is production; GPT-5-mini is a
        free, at-scale cross-check, not an equal second labeler
\end{itemize}
\end{frame}

% ─────────────────────────────────────────────────────────────────────────────
\begin{frame}{Validation Results}
\resizebox{0.95\textwidth}{!}{%
\begin{tabular}{lccc}
\toprule
Field & Claude vs.\ human & GPT vs.\ human & Verdict \\
\midrule
severity\_tier (collapsed) & 96\% & 88\% & Claude preferred \\
root\_cause\_type & 90\% & 70\% & Claude preferred \\
scope & 90\% & 94\% & comparable \\
violation\_category & 84\% & 86\% & comparable \\
\bottomrule
\end{tabular}%
}
\vspace{10pt}
Cross-model agreement at full scale (n=1,067): binary flags \highlight{92--99\%};
severity\_tier collapsed \highlight{94.5\%} (kappa 0.82); root\_cause\_type only
58\% (Claude's own accuracy vs.\ human is fine at 90\%, it is specifically GPT's
weak field).
\vspace{6pt}
\begin{alertblock}{Status}
\small Extraction, validation, and model selection are done. Not revisited unless new
data arrives.
\end{alertblock}
\end{frame}

% ═══════════════════════════════════════════════════════════════════════════
\section{Part 2: Is the Text Useful?}

\begin{frame}{Part 2: Is the Text Useful, and a Problem We Found}
\begin{itemize}
  \setlength{\itemsep}{10pt}
  \item INFORMS (July): text alone predicted future patient harm in VAI-only
        facilities, AUC \highlight{0.656}, $p{=}0.046$
  \item Rerun on the current, validated pipeline: AUC \highlight{0.591--0.610},
        no longer significant ($p{=}0.19$, $p{=}0.08$)
  \item Investigated why -- found something bigger than a weaker replication
\end{itemize}
\end{frame}

% ─────────────────────────────────────────────────────────────────────────────
\begin{frame}{The Outcome Variable Was Confounded}
\begin{columns}[T]
\column{0.55\textwidth}
The AE outcome (above-median AE count, pooled across all facilities and drugs)
was mostly measuring \highlight{which facility this is}, not whether quality
changed.
\vspace{8pt}
\begin{block}{Variance decomposition}
\small \textbf{88.6\%} of the variance in raw AE count is BETWEEN facilities
(size/drug identity). Only \textbf{11.4\%} is WITHIN a facility over time.
\end{block}
\column{0.42\textwidth}
GroupKFold CV correctly refuses to let a model exploit facility identity --
which is most of what the old label contained.
\end{columns}
\end{frame}

% ─────────────────────────────────────────────────────────────────────────────
\begin{frame}{The Fix: Compare Each Facility to Itself}
New outcome: did this facility's own AE count rise after inspection vs.\ its
own pre-inspection window, same facility, same window length.
\vspace{8pt}
\resizebox{0.95\textwidth}{!}{%
\begin{tabular}{lcc}
\toprule
Config (text only) & Old outcome & New outcome (relative) \\
\midrule
Logistic Regression & 0.483 (n.s.) & \highlight{0.589, $p{=}0.048$} \\
Random Forest & 0.532 (n.s.) & \highlight{0.593, $p{=}0.041$} \\
\midrule
OAI/VAI flag alone & 0.537 (n.s.) & \highlight{0.457} (below random) \\
\bottomrule
\end{tabular}%
}
\vspace{10pt}
Two things changed: text now predicts across the \highlight{full sample}, not
a small subgroup; FDA's own classification alone predicts
\highlight{nothing}.
\end{frame}

% ─────────────────────────────────────────────────────────────────────────────
\begin{frame}{Robustness Checks}
\begin{itemize}
  \setlength{\itemsep}{9pt}
  \item \textbf{Feature convergence:} Lab Controls signal and
        data\_integrity\_llm\_share are top-8 in three independently built models
        (facility-year AE, raw correlation, the new significant model)
  \item \textbf{Calendar-time check:} no reporting-lag artifact -- FAERS volume
        stable through 2026 Q1
  \item \textbf{Facility trend plots:} some facilities show a textbook pattern
        (severity spikes, AE count triples next quarter); others are pure noise --
        exactly the mix an AUC of 0.55--0.60 represents
  \item \textbf{AE-linkage precision matters:} the significant result needs
        ANDA-matched AE data; broader drug-level matching is noisier and not
        significant
\end{itemize}
\end{frame}

% ─────────────────────────────────────────────────────────────────────────────
\begin{frame}{Does OAI Calm the Trend More Than VAI?}
Expectation: after an inspection, the AE rise should reverse with a lag,
especially if FDA assigns OAI (forces correction). Tested directly on the
corrected panel, split by \emph{this specific inspection's} classification, not
by whether the facility ever had an OAI (a different, coarser cut used
elsewhere in this deck).
\vspace{8pt}
\resizebox{0.9\textwidth}{!}{%
\begin{tabular}{lccc}
\toprule
Classification & $n$ & Pre$\to$post AE ratio & Direction \\
\midrule
OAI & 29 & \highlight{0.995} & flat / very slightly down \\
VAI & 94 & \highlight{1.183} & up 18\% \\
\bottomrule
\end{tabular}%
}
\vspace{10pt}
\begin{alertblock}{Read this carefully}
\small Direction matches the hypothesis: OAI flat, VAI still rising. But
\highlight{not statistically significant} -- Mann-Whitney $p{=}0.26$, Welch
$t$-test $p{=}0.26$, Fisher exact on simple rise/no-rise $p{=}0.66$. Only 29
OAI and 94 VAI inspections in the corrected (no-zero-fill) panel. Directionally
consistent, not confirmed -- a hypothesis to keep testing as the sample grows,
not a result to cite yet.
\end{alertblock}
\end{frame}

% ─────────────────────────────────────────────────────────────────────────────
\begin{frame}{The Dependent Variable Still Needs to Get Better}
\begin{itemize}
  \setlength{\itemsep}{10pt}
  \item Even fixed, FAERS AE counts top out around AUC 0.55--0.60 -- real signal,
        not a strong one
  \item \tech{MarketScan} (OSU team) offers a stronger design: patient-level
        self-control (same patient, before/after a manufacturer switch), not just
        facility-level
  \item Not usable yet -- only reference documentation exists in the project
        folder so far
  \item Formal data request drafted: exact fields, formulas, time-window caveats,
        which of our 14 drugs the design covers (9 of 14; hospital-administered
        drugs do not qualify)
\end{itemize}
\end{frame}

% ═══════════════════════════════════════════════════════════════════════════
\section{Part 3: Prediction Models}

\begin{frame}{Part 3: Prediction Models -- Recall, AE, Shortage}
\begin{itemize}
  \setlength{\itemsep}{10pt}
  \item Same design question as Part 2, one level up: facility-year models for
        three outcomes
  \item Fixed the same problem here too: models were restricted to
        text-covered facilities only, discarding real inspection history for
        facilities without text yet
  \item Every model now runs two populations independently: a
        \textbf{baseline} (inspection + structural only, full facility
        universe) and a \textbf{with-text} model (text-covered subset)
\end{itemize}
\end{frame}

% ─────────────────────────────────────────────────────────────────────────────
\begin{frame}{Variables Used in Every Model}
{\small
\begin{itemize}
  \setlength{\itemsep}{6pt}
  \item \tech{Inspection (4):} n\_oai\_cumul, n\_vai\_t, n\_inspections\_t,
        n\_warning\_letters\_t
  \item \tech{Structural (2):} parenteral\_ever, n\_feis\_drug (supply
        concentration -- FEIs making the same drug)
  \item \tech{Text/LLM (12):} severity, scope, root cause, contamination,
        data integrity, investigation, repeat findings, Lab Controls share,
        Quality System share, remediation strength
\end{itemize}
}
\vspace{10pt}
\begin{block}{The two models}
\small \textbf{Baseline} = Inspection + Structural only (6 features), every
facility with real data, no text required.
\textbf{With-text} = all 18 features, only facilities with a real as-of-year
483 snapshot.
\end{block}
\end{frame}

% ─────────────────────────────────────────────────────────────────────────────
\begin{frame}{The One Clean Result: Shortage}
\resizebox{0.85\textwidth}{!}{%
\begin{tabular}{lcc}
\toprule
Shortage (m19) & AUC & $p$ vs.\ 0.5 \\
\midrule
Baseline (inspection + structural, no text) & \highlight{0.54--0.61} & \highlight{0.003 / 0.011} \\
With-text (restricted subset) & 0.49--0.50 & n.s. \\
\bottomrule
\end{tabular}%
}
\vspace{10pt}
\highlight{The strongest, cleanest result in either pipeline, and it uses no
text at all.}
\vspace{12pt}
\begin{block}{Recall and AE (facility-year) were also tried}
\small Recall (m14): too few events to model at all (20 total). AE (m17,
annual): baseline improved after the outcome fix but the with-text model got
worse, likely the same small-sample overfitting problem as the shortage
with-text model. Neither gives a clean result at current sample size -- not
shown as a peer comparison to shortage.
\end{block}
\end{frame}

% ─────────────────────────────────────────────────────────────────────────────
\begin{frame}{What the Shortage Result Means}
\begin{columns}[T]
\column{0.5\textwidth}
{\small
\begin{tabular}{lc}
\toprule
Feature & Importance \\
\midrule
n\_feis\_drug & \highlight{62\%} \\
parenteral\_ever & 12\% \\
n\_oai\_cumul & 11\% \\
n\_inspections\_t & 8\% \\
n\_vai\_t & 6\% \\
\bottomrule
\end{tabular}
}
\column{0.47\textwidth}
{\small Random Forest importance, baseline model.}
\end{columns}
\vspace{10pt}
\begin{itemize}
  \setlength{\itemsep}{7pt}
  \item \highlight{Mostly supply concentration}, not quality: n\_feis\_drug
        (how many other facilities make the same drug) alone is 62\% of the
        signal -- sole-source facilities are the risk
  \item Inspection-based features (OAI/VAI/inspection count) still contribute
        \highlight{about 25\% combined} -- a real, if secondary, quality signal
\end{itemize}
\vspace{8pt}
\begin{alertblock}{This does not question whether text signals are useful}
\small Shortage and patient harm (Part 2) are different questions. Shortage is
mostly a market-structure outcome -- how many alternative suppliers exist --
which text was never expected to predict well on its own. Part 2's
significant AE result stands on its own evidence; this slide is not a
rebuttal of it, it's a separate finding about a separate outcome.
\end{alertblock}
\end{frame}

% ═══════════════════════════════════════════════════════════════════════════
\section{Where This Points Next}

\begin{frame}{Where This Points Next}
\begin{itemize}
  \setlength{\itemsep}{9pt}
  \item \textbf{Send the MarketScan request} -- get a real, patient-controlled
        outcome to validate against
  \item \textbf{Stress-test the shortage baseline} -- repeated cohorts, closer
        look at what is driving it
  \item \textbf{Resolve the m17 with-text puzzle} -- why the corrected outcome
        helped the baseline but hurt the smaller with-text subset
  \item \textbf{Keep growing the sample} -- VAI-only and shortage cohorts are
        still small; more calendar time and 483-text coverage is the binding
        constraint everywhere in this deck
\end{itemize}
\end{frame}

% ─────────────────────────────────────────────────────────────────────────────
\begin{frame}{Summary}
\begin{itemize}
  \setlength{\itemsep}{9pt}
  \item \textbf{Closed:} extraction pipeline, human validation, model choice
  \item \textbf{Fixed:} the outcome-variable confound, in both the AE and
        facility-year pipelines
  \item \textbf{Real, if modest:} text signals predict AE change across the full
        sample once the outcome is fixed ($p<0.05$)
  \item \textbf{Best current result:} plain inspection history predicts
        shortage, significantly, without text
  \item \textbf{Open:} a better dependent variable (MarketScan), more stress-testing,
        and more calendar time
\end{itemize}
\end{frame}

\end{document}
